\section*{Notación y dominios de las construcciones}
\phantomsection
\label{sec:ii-notacion}
\addcontentsline{toc}{section}{Notación y dominios de las construcciones}

La tabla identifica los principales símbolos por su dominio y por la
definición que los introduce. Los subíndices, la tipografía y el contexto
operatorio forman parte de la notación. Las referencias remiten a las
definiciones efectivas; las apariciones anticipatorias conservan ese
mismo significado.

\begingroup
\small
\setlength{\tabcolsep}{4pt}
\renewcommand{\arraystretch}{1.12}
\begin{longtable}{@{}>{\raggedright\arraybackslash}p{0.19\linewidth}
                    >{\raggedright\arraybackslash}p{0.49\linewidth}
                    >{\raggedright\arraybackslash}p{0.27\linewidth}@{}}
\toprule
Símbolo & Objeto y significado & Definición y localización\\
\midrule
\endfirsthead
\toprule
Símbolo & Objeto y significado & Definición y localización\\
\midrule
\endhead
\bottomrule
\endfoot

\(\mathcal U_t\)
& Transición sobre el estado con memoria: selección, transporte y
  actualización de registros, en ese orden.
& \eqref{eq:actualizacion}.\\[3pt]

\(\mathcal R_{12}\), \(R_{\rm sgn}\)
& Registro de doce ventanas con ruta, orientación, acarreo e incidencias;
  lector firmado con codominio \(\ZZ^{12}\).
& \S\ref{subsec:k-registro-integral};
  composición \eqref{eq:k-lector-firmado}.\\[3pt]

\(A_m,C_m,V_m\)
& Tres recuentos enteros por ventana: visitas residuales, diferencia de
  trazas orientadas e incidencia firmada de frontera.
& \S\ref{subsec:registro-evaluacion-incidencial},
  definición de los recuentos.\\[3pt]

\(z_m\)
& Bloque decimal incidencial en \(\{0,\ldots,999\}\):
  \(100[A_m]_{10}+10[C_m]_{10}+[V_m]_{10}\).
& \eqref{eq:registro-bloque-incidencial}.\\[3pt]

\(U\), \(U^{(r)}\)
& Vector entero firmado de doce coordenadas y sus tres bloques
  armónicos de cuatro coordenadas.
& \eqref{eq:k-bloques-firmados} y
  \eqref{eq:k-u-firmado}.\\[3pt]

\(K\), \(K_j\)
& Registro integral \(K=\mathcal H_{12}U/4\), con fase y orientación
  fijadas; prolongación de sus coordenadas por \(K_{j+12}=K_j\).
& \S\ref{subsec:k-hadamard}, definición de la transformación integral;
  \S\ref{subsec:alpha-dominio}.\\[3pt]

\(\mathcal K_{\rm ph}\)
& Núcleo de transferencia sobre \(\mathcal U_6\times C_9\),
  denominado también \(K_{\rm ph}\) al describir la representación con
  memoria reducida.
& \S\ref{ii:tpk-arquitectura}, «Transferencia reducida»;
  retorno \eqref{ii:tpk-renovacion}.\\[3pt]

\(A_j\), \(c_j\)
& Bloque de \(\alpha\) en base \(1000\) y cociente entero de la
  división euclídea; el índice \(j\) recorre la expansión posicional.
& \eqref{eq:alpha-recurrencia}.\\[3pt]

\(\mathbf z_m\), \(y_m\)
& Memoria vectorial acumulada en \(\RR^{12}\) y coordenada en el
  marco móvil: \(y_m=S^{-m}\mathbf z_m\).
& \S\ref{sec:eta}, proposición del incremento;
  actualización \eqref{eq:eta}.\\[3pt]

\(S\), \(R=S^{-1}\)
& Permutación cíclica de la base de memoria,
  \(S\mathbf e_m=\mathbf e_{m+1}\); \(R\) designa su inversa
  dentro del desarrollo del resolvente.
& \S\ref{sec:eta};
  \S\ref{subsec:ii-memoria-resolvente}, párrafo inicial.\\[3pt]

\(\eta_m\)
& Incremento vectorial producido por el evento orientado:
  \(\eta_m=a_mS^{-1}\mathbf e_0\).
& \eqref{eq:events} y \eqref{eq:eta}.\\[3pt]

\(U:X\to X\), \(A\)
& En la reducción de memoria, \(U\) es la actualización del estado;
  \(A:V_0\to V_0\) es el bloque retenido de su extensión lineal
  \(T\delta_x=\delta_{U(x)}\).
& \S\ref{subsec:memoria-resolvente}, «Eliminación de memoria»;
  proposición~\ref{mem:prop-schur}.\\[3pt]

\(U:\mathcal H\to\mathcal H\)
& Realización isométrica del transporte de una vuelta,
  \(U^*U=I\), en la compresión observable.
& \S\ref{sec:extension}, «Compresión observable»;
  \eqref{eq:conservacion-ortogonal}.\\[3pt]

\(R=\ZZ/9\ZZ\)
& Anillo residual de APP, con ideal máximo
  \(\mathfrak m=(3)\).
& \S\ref{sec:nucleo}, «Reducción ternaria y radical multiplicativo».\\[3pt]

\(R\), \(\eta\) locales
& En el bloque central \(B_c=7I+2R\), \(R\) intercambia las dos
  coordenadas y \(R^2=I\). La rapidez de su normalización es
  \(\eta=\tfrac12\log(9/5)\).
& \S\ref{sec:extension}, «Desarrollo local»:
  \(B_c/\sqrt{45}=\exp(\eta R)\).\\[3pt]

\(\Sigma_H\)
& Lector adimensional de norma y autoescala
  \(\Sigma_H=\varphi\alpha^{16}\); su valoración decimal es
  \(\operatorname{ord}_{10}\Sigma_H=-34\).
& \eqref{eq:el-lector-determinantal} y \eqref{eq:el-decada}.\\[3pt]

\pagebreak[4]

\(\mathcal L_S\), \(\mathcal U_S\)
& Recta orientada de acción y base positiva en la que se expresan
  sus secciones.
& \S\ref{sec:action}, «Secciones de acción y cociente adimensional
  de retorno».\\[3pt]

\(H_5\), \(\eta_{\rm ret}\)
& Coeficientes adimensionales de acción anterior y posterior al retorno;
  \(\eta_{\rm ret}=H_5-(90/\pi)\mathcal C_\pi(\alpha)\).
& \eqref{eq:el-h5}--\eqref{eq:el-eta}.\\[3pt]

\(\hbar_{\rm pre}\), \(\hbar_{\rm ret}\), \(h_a\)
& Secciones \(H_5\,10^{-34}\mathcal U_S\) y
  \(\eta_{\rm ret}\,10^{-34}\mathcal U_S\).
  Para \(a\in\{\mathrm{pre},\mathrm{ret}\}\),
  \(h_a=2\pi\hbar_a\) expresa la acción por vuelta.
& \eqref{eq:el-secciones-accion};
  \S\ref{sec:action}, paso a la acción por vuelta.\\[3pt]

\(R_{\rm act}\), \(\delta_{\rm ret}\)
& Razón multiplicativa \(H_5/\eta_{\rm ret}\) y diferencia
  aditiva \(H_5-\eta_{\rm ret}\) de las coordenadas de retorno.
& \eqref{eq:el-ratio}.\\[3pt]

\(A_{\deg}\), \(C^*_{\deg}\)
& Coordenadas angulares \(1000\alpha\) y
  \(2(\eta_{\rm ret}+\alpha)\). Sus coordenadas en radianes son
  \(x=\pi A_{\deg}/180\), \(y=\pi C^*_{\deg}/180\).
& \eqref{eq:ii-par-angular};
  \S\ref{sec:ii-angulos}, definición de las coordenatizaciones.\\[3pt]

\(\gamma_{\HMT}\), \(\gamma\)
& Evaluación conjugada de la cadena orientada \(c_{12}^{+}\) en
  \(q_\pm\), con componente angular \(\pi AC^*/180\), para
  \(0<q_\pm<1\). El incremento de área positivo utiliza
  \(|\gamma_{\HMT}|\).
& \S\ref{sec:ii-definicion-barbero};
  evaluación \eqref{eq:ii-definicion-barbero-evaluacion}.\\[3pt]

\(\eta_{\rm el}\), \(a_S,b_S\)
& Parámetro de forma \(\operatorname{artanh}(C^*/A)\) y
  semiejes canónicos equilibrados de la elipse, con
  \(a_Sb_S=2\hbar\).
& \S\ref{sec:ellipse}, definición inicial.\\[3pt]

\(R_{\rm el}\), \(\tau_{\rm el}\)
& Radio modular adimensional \(\sqrt{b_S/a_S}\) y módulo
  \(\tau_{\rm el}=iR_{\rm el}^2\).
& \S\ref{sec:ellipse}, reciprocidad de la elipse.\\[3pt]

\(K_a^i\)
& Curvatura extrínseca en la convención normal de la realización
  geométrica; interviene en
  \(\mathcal A_a^i=\Gamma_a^i(E)+\gamma K_a^i\).
& \S\ref{sec:ii-costura-barbero}, definición de la conexión.\\[3pt]

\(L_*\), \(L\), \(\ell_P\)
& Referencia longitudinal positiva y longitud
  \(L=(5759/23040)\alpha_{\HMT}^{16}L_*\).
  La realización R1--R8 determina \(\ell_P^2=L^2\);
  la evaluación SI declara \(L_*=1\,\mathrm m\).
& \eqref{g26:eq:regla-longitudinal-es} y
  \eqref{g26:eq:g-rigido-es}.\\[3pt]

\(R_X\), \(\Lambda_X\), \(E_L\)
& Lectura radial positiva, longitud conjugada operatoria y escala
  energética \(E_L=\hbar_{\rm ret}c_{\rm int}/L\).
  Satisfacen \(R_X\Lambda_X=L^2I\).
& \eqref{g26:eq:lecturas-conjugadas-es} y
  \eqref{g26:eq:productos-radiales-es}.\\[3pt]

\(R_{108}\), \(R\) de horizonte
& Radio circular \(c_{\rm int}/\omega_{108}\) del período
  dodecafásico, igual a \(L\) para el reloj conjugado;
  \(R>0\) es la variable radial de la realización de
  horizonte.
& \S\ref{sec:gravity}, definición inicial;
  \eqref{eq:ii-horizonte-normalizado}.\\[3pt]

\(G_a\)
& Sección gravitatoria recíproca de una acción \(\hbar_a>0\):
  \(G_a=c_{\rm int}^{3}L^{2}/\hbar_a\), con \(R_{108}=L\),
  \(c_{\rm int}>0\) y coincidencia de radios en la convención
  \(r_g=Gm/c_{\rm int}^{2}\).
& \S\ref{sec:ii-definicion-gravedad};
  \eqref{eq:ii-definicion-gravedad-seccion} y
  \eqref{eq:ii-definicion-gravedad-area}.\\[3pt]

\(\mathbb G_{5,a}\)
& Operador \(G_aP^{\mathrm{ico}}_5\) sobre
  \(V_{12}=\RR[A_5/C_5]\). El proyector ortogonal selecciona
  el sumando \(V_5\); la sección \(G_a\) es común a sus cinco
  componentes.
& \eqref{ii:pent:eq:gravity-G5};
  teorema~\ref{ii:pent:thm:gravity-p5-coupling}.\\[3pt]

\(k_B\), \(\Theta_{{\rm clk},a}\)
& Transductor lineal positivo
  \(k_B:\mathcal L_\Theta\to\mathcal L_E\) y sección térmica
  positiva, relacionados por
  \(k_B\Theta_{{\rm clk},a}\log3=h_a/(108t_0)\).
  La sección térmica y las bases declaradas determinan su
  coordenada individual.
& \S\ref{sec:ii-definicion-boltzmann};
  \eqref{eq:ii-kb-aut-par-termico} y
  \eqref{eq:ii-kb-aut-coordenada}.\\

\end{longtable}
\endgroup

En el registro incidencial las ventanas llevan índices \(1,\ldots,12\).
En la actualización de memoria se numeran desde \(0\) hasta \(11\), y
su prolongación utiliza \(m\geq0\). El desplazamiento de índice conserva
la misma partición en ventanas nonádicas. En las fórmulas angulares,
\(A,C^*\) conservan la coordenatización declarada en cada expresión; el cociente
\(C^*/A\) es invariante bajo el cambio común entre grados y radianes.
